% !TeX root = ../main.tex
% ---------------------------------------------------------------------------
% Large Caverns and Complex Geometries
% Scope extension: expands the cavern material touched on in Chapters 6 and 10.
% Key sources: Hoek & Moy, Design of large powerhouse caverns; Barton et al.
% on Gjovik; Benson on underground powerhouses; Khodr et al. (2023) Snowy 2.0.
% ---------------------------------------------------------------------------
\chapter{Large Caverns and Complex Geometries}
\label{ch:caverns}


\section{Applications and cavern layout}
\label{sec:applications-and-cavern-layout}

\subsection{Applications of large underground openings}
\label{sec:cavern-applications}
\scopenote{Hydropower machine halls, metro station caverns, storage caverns
for hydrocarbons and compressed air, sports and civil defence facilities, and
deep geological repositories.}

\subsection{Cavern shape and stress optimisation}
\label{sec:cavern-shape}
\scopenote{Mushroom, horseshoe and elliptical profiles; matching the opening
to the stress ellipse to minimise tensile zones and stress concentration. The
convergence comparison presented in Chapter~\ref{ch:construction}.}

\section{Orientation and span}
\label{sec:orientation-and-span}

\subsection{Orientation relative to stress and structure}
\label{sec:cavern-orientation}
\scopenote{Aligning the long axis with the major horizontal stress, and
against the dominant joint set; the frequent conflict between these two
objectives and how it is resolved.}

\subsection{Span, stand-up and empirical guidance}
\label{sec:cavern-span}
\scopenote{Extending the $Q$-system and equivalent dimension to large spans,
the ESR values appropriate to permanent caverns, and the limits of empirical
extrapolation beyond the case base.}

\section{Excavation sequencing and roof behaviour}
\label{sec:excavation-sequencing-and-roof-behaviour}

\subsection{Excavation sequencing}
\label{sec:cavern-sequencing}
\scopenote{Top heading and benching, multiple drifts, sidewall galleries and
pilot tunnels; the sequence as a means of controlling deformation and
progressively confirming ground behaviour.}

\subsection{Roof arch behaviour}
\label{sec:roof-arch}
\scopenote{Arch action, voussoir beam behaviour in bedded roofs, and buckling
or snap-through failure of a jointed roof beam. Related script:
\code{arc-beam.py}.}

\section{Structural stability and reinforcement}
\label{sec:structural-stability-and-reinforcement}

\subsection{Wedge analysis at cavern scale}
\label{sec:cavern-wedges}
\scopenote{Applying the projection methods of
Chapter~\ref{ch:structural-analysis} at large span: the cubic growth of wedge
weight with span, the multiple differently oriented surfaces a cavern
presents, and probabilistic joint orientation input.}

\subsection{Cable bolting and high-capacity reinforcement}
\label{sec:cable-bolting}
\scopenote{Long tendons, pre-tensioning, spacing and length rules for large
spans, and the interaction between systematic bolting and discrete cable
support. Detail in Chapter~\ref{ch:rock-reinforcement}.}

\subsection{Sidewall stability and high walls}
\label{sec:cavern-sidewalls}
\scopenote{Tall sidewalls in high horizontal stress, stress relief, buckling
of steeply dipping foliation, and support of the wall as distinct from the
crown.}

\section{Intersections and observational control}
\label{sec:intersections-and-observational-control}

\subsection{Intersections, junctions and bifurcations}
\label{sec:junctions}
\scopenote{Three-dimensional stress redistribution at openings, reinforcement
of the intersection, and why these locations require 3D analysis
(Chapter~\ref{ch:advanced-numerical}).}

\subsection{Observational control of cavern excavation}
\label{sec:cavern-observation}
\scopenote{Instrumented sequential excavation, trigger levels tied to
excavation stages, and design modification during construction.}

\section{Case study: the Snowy 2.0 power station complex}
\label{sec:case-snowy-complex}
\scopenote{Cavern shape selection under high stress, structurally controlled
wedge assessment, and the support system adopted.}
